\section{Introduction}
\label{sec:introduction}

\subsection{The problem, not the product}
\label{sec:intro-problem}

\paragraph{Where this comes from.} \Flux{} is not a proposal. The chain has produced blocks
continuously since \textbf{2018}: mainnet genesis carries UNIX timestamp \num{1516980000} ---
2018-01-26T15:20:00Z --- and its coinbase still reads \texttt{Zelcash06f40b01\ldots}, the name the
project launched under and the rebrand the consensus code has never shed, as \texttt{determ\_zelnodes},
\texttt{zelid} and \texttt{ZelBack} attest throughout this paper \shipped
\srccode{fluxd/\allowbreak{}src/\allowbreak{}chainparams.cpp}{65,185}. The tiered node network was
announced in October 2018 \citep{zelnodes2018}; its orchestration layer --- \FluxOS{}, then ZelFlux
--- debuted in the autumn of \textbf{2019} \citep{zelflux2019}. Eight years of continuous operation
is the reason this paper can describe half its subject in the past tense at all, and the reason its
claims about deployed behaviour are checkable against running code rather than argued from a design.

That history has been written about before, and this paper does not restate it. The chain was
specified in the \textbf{ZelCash whitepaper v2} of July 2018 \citep{zelcashwp2018}. Two documents are
this paper's direct predecessors, and it supersedes both rather than amending them. The
\textbf{Proof-of-Useful-Work v2 whitepaper} \citep{fluxpouw2} --- still linked from \fluxd's own
README --- specified the consensus mechanism now called Proof of Node, and is the one this document
builds on most directly rather than departs from: \S\ref{sec:consensus} carries its construction
forward as deployed code and measures it, and \S\ref{sec:parallel-assets} reconciles its supply
model to within one block of accrual (\Cref{tab:pa-depletion-epochs}). Where the two differ, it is
because something has since been built, measured or decided, and this document says so at the point
of difference. \textbf{Flux Whitepaper 2.1} \citep{fluxwp21} is the ecosystem description it
replaces. The author's own contemporaneous descriptions of node tiering \citep{kmentazelnode2020} and
of the \FluxOS{} architecture \citep{kmentaflux2020} are prior art for parts of \S\ref{sec:consensus}
and Part~III, and the per-tier benchmark thresholds \S\ref{sec:consensus} reads out of \fluxbench{}
were published as an operator guide \citep{zelbenchguide}.

\notestablished{the publication dates of the historical Medium articles cited above, with one
exception. Medium returns HTTP 403 to scripted retrieval, so only \citet{kmentaflux2020} carries an
independently confirmed timestamp (2020-03-18T03:17:28Z, read from the \texttt{datePublished}
metadata of an Internet Archive capture of 2020-11-06); the remaining dates are those of the
index compiled for this survey. Both articles under the author's own byline do carry confirmed
archival fixity --- the Wayback availability API reported status-200 captures for each on
2026-09-06 --- so the \emph{content} is anchored even where the publication date is not. None of the
arguments in this paper depends on those dates in any case: the 2018 and 2019 anchors above rest on
the genesis block and on \fluxd's source, not on any article.}


A cloud provider today sells capacity to people. A human opens an account, attaches a card or signs
a contract, and later closes the account when the capacity is no longer needed. Every part of that
sentence is an assumption about who the customer is, and the assumption is load-bearing: the account
presumes an identity a provider can bill, the card presumes a payment network built around recurring
consent from a person, and the contract presumes a party who can be sued, renegotiated with, or
simply asked. This paper's premise is that if the dominant consumer of cloud capacity becomes
autonomous software rather than people, each of those three assumptions has to be replaced with
something a program can satisfy entirely on its own — an identity a piece of code can hold, a payment
a piece of code can make without a human tapping a button, and an authority a piece of code can prove
it has without a lawyer's signature. That is a bet about demand, stated here as the premise the rest
of the paper works from, and not a prediction this paper defends. \Flux's own public roadmap states
the destination directly: ``a cloud that serves software agents as first-class customers\ldots no
account, no card, no human in the loop'' \vision \srcroadmap{flux-roadmap-public.md:7}. Whether that
demand materializes is outside this paper's scope; whether the architecture that follows from betting
on it is coherent, and how much of it is actually built, is the whole of it.

Declining to forecast is not the same as declining to say why the bet is reasonable, and a reader
entitled to skip Part~IV on the grounds that its premise is unargued deserves the argument. Three
things make it worth taking seriously. \emph{First, the direction is already observable rather than
speculative}: software that calls APIs, provisions resources and completes multi-step tasks without a
human approving each step is a shipping product category, not a projection, and every such call is
one a card-and-account model was not designed to price. \emph{Second, the bet does not require
agents to dominate.} It requires only that the \emph{marginal} new customer is increasingly
machine-initiated, because the three assumptions being replaced --- an account someone opens, a card
someone consents to, a contract someone signs --- fail at the first such customer, not at the
majority. A provider that cannot serve one autonomous buyer cannot serve any. \emph{Third, the
failure mode is bounded}: every mechanism Part~IV specifies --- delegated authority with a budget,
payment without a custodian, deployment without an account --- is independently useful to a human
tenant who wants to bound what a script may spend or deploy without a hosting relationship. If the
demand does not arrive, the architecture is over-built for its market rather than wrong about its
mechanism.

That is the case, and it is deliberately not stronger than the evidence. \Cref{sec:evaluation} names
what would falsify it, and \Cref{sec:conclusion} states plainly that if agents are not the customer,
much of the target architecture solves a problem the network does not have.

There is a second, narrower observation about \emph{what} such software would run, and it turns out to
have architectural consequences rather than only commercial ones. Much of the work an autonomous agent
does for an ordinary person --- drafting text, triaging mail, summarising a document, keeping a
schedule --- is served adequately by small models specialised to a task, not by frontier ones. That
distinction matters here because the two have different physics. Frontier inference splits a model
across many accelerators and exchanges state between them on every token, so it is bound to an
interconnect one to two orders of magnitude faster than wide-area networking; that is why it
concentrates in single datacentres, and it is a topology constraint rather than a preference. A model
that fits on one device exercises no such interconnect, and concurrent requests to it are mutually
independent. \textbf{For that class of work the structural advantage of centralisation is not reduced
--- it is absent}, and a dispersed fleet of collateralised commodity machines becomes an appropriate
substrate rather than a compromised one. \S\ref{sec:ac-workload-fit} argues this properly, including
the part that cuts the other way: \Flux allocates accelerators whole, and serving many small models
economically requires subdividing them, which is the one thing the current allocation model cannot do.

\subsection{\Flux today, measured}
\label{sec:intro-today}

Before any forward-looking claim, the paper establishes what is running. \Flux is not a proposal: on
2026-09-03 the network carried \num{6489} collateralised nodes across three tiers — Cumulus,
Nimbus and Stratus — backing \num{88514500}\flux\ of locked collateral
\shipped \srcmeasured{api.runonflux.io/\allowbreak{}daemon/\allowbreak{}listzelnodes}{2026-09-03}; \num{1195} applications were
deployed across \num{7486} requested instances
\shipped \srcmeasured{api.runonflux.io/\allowbreak{}apps/\allowbreak{}globalappsspecifications}{2026-09-03}; and since block
height \num{2020000} the chain has produced every block through a collateral-gated lottery over the
registered node set rather than through Equihash proof-of-work
\shipped \srccode{fluxd/\allowbreak{}src/\allowbreak{}chainparams.cpp}{153} (\S\ref{sec:consensus}). These numbers are the
yardstick for everything that follows: a paper about restructuring a chain's incentives and execution
layer around autonomous demand is a different exercise when the chain in question already has real
operators, real collateral and real deployed workloads than when it does not. \S3 gives the full
measured baseline, with every figure in this section reproducible from the same public API on the
same date.

\subsection{The three changes}
\label{sec:intro-changes}

\Flux describes itself as restructuring three things at once: how a node operator is paid, what a
node operator's hardware can be trusted to have actually run, and how value moves between a tenant
and the network. Each is at a different point between specification and deployment, and this paper
states that difference for each rather than presenting all three as equally real.

\paragraph{Emission replaced by demand-funded operator compensation.}
Every payment a tenant makes for an application today lands, in full, at a single transparent
address, and every unit a node operator earns is newly issued rather than routed from that payment
\shipped \srcmeasured{api.runonflux.io/\allowbreak{}apps/\allowbreak{}deploymentinformation}{2026-09-03}. Progressive Node
Rewards (\PNR) is the specified replacement: a consensus-enforced fraction of each application
payment would be diverted into a protocol-held fee pool and paid out to operators through the
existing block-production rotation, one small increment per block, so that operator income begins to
track what tenants actually pay rather than a fixed issuance schedule alone. \PNR is specified to
implementation depth in \S\ref{sec:pnr} — the consensus rule, the pool's per-block update, and two
fairness proofs — but nothing in it exists in any repository surveyed for this paper
\planned \srcintent{design intake}. It is a design, stated precisely enough to build, not
a description of running code.

\paragraph{The substrate hardened by attested execution.}
\ArcaneOS pairs a measured Secure Boot chain with disk encryption and a signed, \texttt{dm-verity}
root filesystem, and a local attestation service (the \SAS) that other components can query for a
signed claim about what a node has booted. That boot chain and that service run today
\shipped (\S\ref{sec:attested-execution}). What they do not yet provide is named precisely in that
section rather than left implicit: the resulting signature is not bound to the specific machine that
produced it, so a bond placed against it would secure a claim any operator's own software could
reproduce, not a claim about a particular piece of hardware. The bonded attestation network proposed
to close that gap — committee sampling, a challenge protocol, slashing against a separate bond — is
\planned and, as \S\ref{sec:feasibility} argues at length, is a precondition for the first change
above rather than an independent workstream, because a compensation mechanism that pays every
eligible node raises the payoff to passing an eligibility gate that is not yet machine-bound by
exactly what it pays.

\paragraph{Value settled transparently, on rails whose privacy layer is being withdrawn.}
\fluxd is a Zcash fork, and its shielded-transaction machinery — Sprout and Sapling notes, encrypted
outputs, a note commitment tree, Groth16 proofs — is inherited, compiled into every node, and running
today \shipped (\S\ref{sec:shielded-value}). The shielded pools presently hold
\num{96479.94}\flux, against a total issued supply of \num{429466823.97}\flux\ at height
\num{2912015} — \SI{0.0225}{\percent} of supply
\shipped \srcmeasured{api.runonflux.io/\allowbreak{}daemon/\allowbreak{}getblockchaininfo}{2026-09-03}. The pool is closed by
consensus: transparent-to-shielded transfers have been rejected since height \num{835554}
\shipped \srccode{fluxd/\allowbreak{}src/\allowbreak{}main.cpp}{1398--1408}, so every shielded unit on the chain today predates
the measured tip by over two million blocks. \textbf{Both pools are to be retired}
\srcintent{design intake} (\planned). The measurement above is most of the
argument: a mechanism holding \SI{0.0225}{\percent} of supply, unable to grow, and requiring a Rust
proving stack in the consensus path is a poor trade against a stated goal of less code and a more
thorough audit (\S\ref{sec:shielded-retirement}). The decisive argument is narrower and is worth
stating here rather than only in \S20: a shielded pool would not hide what a tenant most plausibly
wants hidden, because a \Flux{} application specification is a public on-chain message and the
control plane serves the global specification set over a public endpoint. The workload is public by
construction; shielding the payment hides the funding source, not the workload graph.

So the honest statement of this pillar is the one the paper makes and defends rather than the one an
earlier draft asserted: \textbf{\Flux{} settles transparently, and treats payment privacy as an
unsolved problem rather than a delivered feature}. \S\ref{sec:private-payment} states the
requirement formally and proves two results that outlast the decision --- that amount hiding is
unsatisfiable under \PNR{}, and that the residual problem is confined to one optional path --- but
its constructions are inapplicable to a chain with no pool, and it says so.

\subsection{What is new in this paper}
\label{sec:intro-new}

Nothing below claims novelty for what \Flux inherited outright — the shielded transaction machinery
is Zcash's (\S\ref{sec:shielded-value}), the UTXO model and the limits on what its scripts can express
are Bitcoin's (\S\ref{sec:agent-identity}), and the sortition-based block production that replaced
mining is deployed, described, and not invented here (\S\ref{sec:consensus}). Four results in this
paper are, to the corpus surveyed, new.

First, a proof that pooling and dripping fee revenue — rather than paying it out to next block's
payee directly — defeats a specific timing attack that deterministic payee rotation otherwise admits.
Under a next-block payout rule, an operator who is also a tenant can time its own payment to land just
before a block it is scheduled to receive, and for a single node in the highest tier this recaptures
up to roughly a third of the operator's own payment at the design's proposed parameters; under a
fee-pool drip the same attack recaptures a fraction bounded by that tier's share of one drip step,
several orders of magnitude smaller for the recommended drip constant
\planned (\S\ref{sec:pnr}).

Second, a measurement, not merely a construction, that collateral-weighted moderation is
Sybil-resistant by construction and, at the distribution measured on 2026-09-03, not
censorship-resistant. Voting weight cannot be increased without locking additional collateral, which
is the resource-cost property the design is built on; measured against the live node-weight
distribution, four owner identities already hold enough collateral-weighted voting power to reach the
proposed \SI{25}{\percent} removal threshold on their own
\shipped \srcmeasured{api.runonflux.io/\allowbreak{}daemon/\allowbreak{}listzelnodes}{2026-09-03} (\S\ref{sec:23-governance}).
Both properties are true of the same mechanism, and the paper states both.

Third, an observation about the consensus header format, offered with an explicit hedge because it
has not been checked against a bridge implementation \srcinferred. Every block produced under the
current consensus rule carries a fixed chainwork value of $2^{40}$, a choice made so that fork choice
is a block-count race rather than a work comparison \shipped \srccode{fluxd/\allowbreak{}src/\allowbreak{}pow.cpp}{284--290}
(\S\ref{sec:consensus}). A header chain with identical work per block carries no signal a light
client could use to prefer one candidate chain over another beyond counting blocks, which is not the
guarantee header-only verification is built to provide; if that reading holds, it forecloses the kind
of header-only light-client bridge design a work-bearing chain would otherwise support, a constraint
this paper argues \S22 must design around rather than treat as incidental.

Fourth, a separation of an agent's \emph{deployment} authority from its \emph{spending} authority
that makes the first tractable with machinery already running and states precisely why the second is
not tractable on the \Flux\ L1 at all without either a new trusted party or new consensus state:
script validation on a Bitcoin-derived chain is a pure function of the spending transaction and the
outputs it consumes, so no script can read how much a key has already spent and therefore no script
can enforce a rate limit or a cumulative cap \planned (\S\ref{sec:agent-identity}).

\subsection{What this paper does not claim}
\label{sec:intro-not}

The paper makes no price claim, sets no target, and uses no investment framing anywhere; it confines
itself to issuance, supply, operator income and demand
\srcintent{design intake}. It does not assert that the agent-demand premise stated in
\S\ref{sec:intro-problem} is correct — only that the architecture described in the rest of the paper
follows from betting that it is. And it does not claim that the work remaining to reach that
architecture is uniformly straightforward. \S\ref{sec:feasibility} classifies every planned mechanism
by what kind of work it needs and finds four that are Research rather than Engineering or Extension:
post-quantum migration, for which nothing exists in the surveyed repositories and no scheme yet
addresses outputs whose keys the protocol cannot move; bounded spending authority for an autonomous
agent on the \Flux\ L1, which the UTXO script model cannot express without a new trusted party or new
consensus state; a shielded payment provably funding one specific, unlinked deployment, which is hard
because an application's address is itself a workload identifier; and the moderation quorum's
concentration problem, where no reweighting, threshold change or eligibility gate the paper considered
raises the number of colluding parties needed to censor an application above single digits at the
measured distribution.

\subsection{How to read this paper}
\label{sec:intro-howtoread}

Every substantive claim in this document carries two tags. A \emph{maturity} tag — \shipped,
\inprogress, \planned or \vision — states how real the claim is: \shipped means it runs today and is
written in the present tense; \planned and \vision describe what does not yet exist and are written
in the future or the subjunctive, never in the present. A \emph{provenance} tag states where the claim
comes from: a code citation (\srccode{repo/\allowbreak{}file}{line}), a roadmap citation
(\srcroadmap{item}), a design-intent citation (\srcintent{\S n}), a measurement
(\srcmeasured{endpoint}{date}), or an explicit hedge (\srcinferred) for a conclusion this paper draws
rather than one the evidence states outright. The rule that makes the pairing meaningful is strict:
\emph{only} a code citation --- to deployed source or to a named branch --- can support a \shipped
claim about behaviour, and \emph{only} a measurement can support a \shipped claim about a fact of the
live network, such as a node count, a balance or a price, which no line of code can establish. A
roadmap item or a design-intent note can
describe an aspiration precisely, and this paper cites them precisely, but neither can promote a claim
to \shipped on its own. At any sentence, this pairing is meant to let a reader answer two questions
without searching elsewhere in the document: does this exist, and how do we know. Where a
\emph{Settled}, \emph{Not established by this survey} or \emph{Proposed by this paper} callout
appears, it marks respectively a decision that has been taken together with the value chosen, a fact
no artefact read for this paper carries, or a construction the author has not ruled on — never a gap
the authors failed to notice.

The document is organized in five parts and a bridge between them. Part~I fixes the system and
adversary model and the measured baseline this section previews. Parts~II and~III describe the
deployed system — consensus, monetary policy, the orchestration layer, networking, attested
execution — at a depth intended to support reimplementation, including the parts that do not work the
way a casual description would suggest. Part~IV specifies the target architecture for agents and
value at the same depth, with every undecided parameter named rather than smoothed over. Part~IV.5
composes the whole of it into one operational walkthrough, run once against the deployed system and
once against the target. Part~V tests the claim that the remaining distance is extension rather than
invention and reports, in \S\ref{sec:feasibility}, that it is not uniformly so.

\subsection{Method}
\label{sec:intro-method}

Three different activities produced the claims in this paper, and keeping them visibly separate is
the single most important discipline it follows. What was \emph{read} is a corpus of \Flux
repositories, including private ones that the team has agreed may be cited freely by file and line in
this public document \srcintent{design intake}; every \shipped claim about behaviour
traces to a specific file and line range in that corpus. What was \emph{measured} is a set of named public
endpoints — chiefly \texttt{api.\allowbreak{}runonflux.\allowbreak{}io}'s node, application and blockchain-info routes — queried
on 2026-09-03 at chain height \num{2917115}, with every figure reproducible by running the
accompanying scripts against the same live API. What was \emph{designed} is different in kind from
both: mechanisms specified in this paper's own design-intake documents
(\srcintent{design intake}) that do not exist in any repository and are not measurements of
anything running, but are this paper's own construction, stated to a level of precision that a reader
could implement and that a reviewer could attack. A reader who cannot tell, at any given sentence,
which of these three produced it has found a defect in this paper, not a subtlety to be inferred.

%% Drain this section's float queue before the next begins.
\FloatBarrier
